\scanpage{266}
\mspage{216}

Je présume que a) sera rarement satisfait sans que \unsure{b|S} le soit (et inversement d'ailleurs), et que la démonstration de a) et b) se fera pratiquement en une seule haleine.

Si $M=\mathrm{Cat}$, $W$ un localiseur fondamental, alors la \deleted{\unsure{validité de a)}} notion adoptée ici de $W$-fibration (lire~: ``ultra forte'', dans la terminologie spécifique développée dans $\mathrm{Cat}$) est à tel point forte, qu'elle implique déjà \inserted{pratiquement} b) -- (moyennant $W(4\unsure{a'|a})$, il est vrai, cf p. 53), et a) n'est guère vérifiable que si on a $W(4a)^{*}$ --
% slanted
\margin{Non, il me semble que $W(4b)$ doit suffire ...}
(loc. cit.). Donc ici c'est ``l'axiome des fibrations'' sur $W$, sous sa forme forte $W(4a)$, qui est la clef de voûte d'une bonne théorie des fibres homotopiques.

Une telle théorie (revenant au cas général) revient à dire, \inserted{pour l'essentiel,} qu'il y a ``assez'' de $W$-fibrations (``ultra fortes'') pour pouvoir décrire \uline{tout} objet d'un $M_{/S}$, à $W_S$-équivalence près, comme un objet $W$-fibrant sur $S$. Dans \deleted{\unsure{tous}} les cas à ma connaissance, on a \unsure{de|bien|alors} plus ceci~: on
% archive sheet number: 266
